\homeworktitle{Домашняя работа}{Рациональные уравнения с модулем и параметром}
\studentline

\begin{routebox}
  \textbf{Ключевой факт:}
  \[
    \frac{A(x)}{B(x)}=0
    \quad\Longleftrightarrow\quad
    \begin{cases}
      A(x)=0,\\
      B(x)\ne0.
    \end{cases}
  \]
  \textbf{Порядок работы:}
  \begin{enumerate}[label=\arabic*.,leftmargin=1.8em,itemsep=0.2em,topsep=0.25em]
    \item найдите нули числителя на каждой ветви модуля;
    \item определите условие, при котором имеется нужное число нулей;
    \item подставьте каждый найденный ноль числителя в знаменатель;
    \item объедините условия и запишите ответ.
  \end{enumerate}
\end{routebox}

\section*{Часть I. Разминка}

\begin{homeworktask}[Ноль дроби]{1}
  Решите уравнение
  \[
    \frac{(x-2)(x+3)}{x-2}=0.
  \]
  Объясните, почему один из нулей числителя не является корнем исходного
  уравнения.
\end{homeworktask}
\answerspace{16mm}

\Needspace{18\baselineskip}
\begin{homeworktask}[Число корней уравнения числителя]{2}
  Определите число корней уравнения
  \[
    2|x|-x=b
  \]
  в зависимости от \(b\).

  \[
    2|x|-x=
    \begin{cases}
      \rule{27mm}{0.4pt}, & x\ge0,\\[2mm]
      \rule{27mm}{0.4pt}, & x<0.
    \end{cases}
  \]

  \begin{center}
    \renewcommand{\arraystretch}{1.4}
    \begin{tabularx}{0.78\textwidth}{>{\centering\arraybackslash}p{0.3\textwidth}>{\centering\arraybackslash}X}
      \toprule
      \rowcolor{TableHeader}
      Значение \(b\) & Количество корней \\
      \midrule
      \(b<0\) & \rule{30mm}{0.4pt} \\
      \(b=0\) & \rule{30mm}{0.4pt} \\
      \(b>0\) & \rule{30mm}{0.4pt} \\
      \bottomrule
    \end{tabularx}
  \end{center}
\end{homeworktask}

\section*{Часть II. Применяем алгоритм}

\begin{homeworktask}[Решение с подробным планом]{3}
  Найдите все значения \(a\), при которых уравнение
  \[
    \frac{|2x|-x-1-a}{x^2-x-a}=0
  \]
  имеет ровно два различных корня.
\end{homeworktask}
\structuredsolutionmap

\begin{homeworktask}[Решение с сокращённым планом]{4}
  Найдите все значения \(a\), при которых уравнение
  \[
    \frac{3|x|-x-2-a}{x^2-x-a}=0
  \]
  имеет ровно два различных корня.
\end{homeworktask}
\reducedsolutionmap
\answerspace{25mm}

\Needspace{18\baselineskip}
\section*{Часть III. Осмысление}

\begin{homeworktask}[Найдите пропущенный этап решения]{5}
  Ученик написал:
  \begin{quote}
    «Числитель уравнения
    \(
      \dfrac{|3x|-2x-2-a}{x^2-2x-a}=0
    \)
    имеет два различных нуля при \(a>-2\). Следовательно, ответ:
    \(a>-2\)».
  \end{quote}

  \begin{enumerate}
    \item На каком этапе решение стало неполным?
    \item Условие \(a>-2\) неверно или лишь недостаточно?
    \item Запишите правильный ответ.
  \end{enumerate}
\end{homeworktask}
\answerspace{20mm}

\begin{homeworktask}[Полная классификация]{6}
  Для исходного уравнения
  \[
    \frac{|3x|-2x-2-a}{x^2-2x-a}=0
  \]
  определите число корней при всех значениях параметра \(a\).

  \begin{center}
    \renewcommand{\arraystretch}{1.45}
    \begin{tabularx}{0.9\textwidth}{>{\centering\arraybackslash}p{0.55\textwidth}>{\centering\arraybackslash}X}
      \toprule
      \rowcolor{TableHeader}
      Значение параметра & Количество корней \\
      \midrule
      \(a<-2\) & \rule{24mm}{0.4pt} \\
      \(a=-2\) & \rule{24mm}{0.4pt} \\
      \(a\in\{-1,0,3,8\}\) & \rule{24mm}{0.4pt} \\
      \(a>-2,\ a\notin\{-1,0,3,8\}\) & \rule{24mm}{0.4pt} \\
      \bottomrule
    \end{tabularx}
  \end{center}
\end{homeworktask}

\newpage
\section*{Часть IV. Перенос метода}

\begin{homeworktask}[Задание со звёздочкой: параметр внутри модуля]{7*}
  Найдите все значения параметра \(a\), при которых уравнение
  \[
    \frac{2|x-a|-x-2}{x^2-1}=0
  \]
  имеет ровно два различных корня.

  \begin{hintbox}
    Рассмотрите случаи \(x\ge a\) и \(x<a\). Здесь параметр перемещает
    точку излома графика вдоль оси \(Ox\), поэтому прежний шаблон нельзя
    применять механически.
  \end{hintbox}
\end{homeworktask}
\solutionmap

\begin{selfcheck}
  \begin{itemize}[label=\(\square\),leftmargin=1.7em,itemsep=0.25em,topsep=0.2em]
    \item Я сохранил условия ветвей модуля.
    \item Я проверил каждый ноль числителя в знаменателе.
    \item Я различаю значения переменной \(x\) и параметра \(a\).
    \item Я постепенно снимал опоры и в задаче 4 построил решение сам.
    \item В ответе указаны промежуток и все исключаемые точки.
  \end{itemize}
\end{selfcheck}
